\section{Introduction}

Artificial intelligence (AI) agents built on large language models (LLMs) are moving beyond code generation into network and service management. Recent work applies them to the operation of beyond-5G and 6G networks~\cite{corici2025agents,xiao2025agentic}, to decentralized optimization of wireless networks~\cite{lee2025decentralized}, and to fault diagnosis in network management systems~\cite{chen2025assistant}. Delegating such tasks to agents can shorten reaction times and remove the human from routine operational loops.

However, this autonomy raises an accountability problem. AI agents do more than produce text: a single task can invoke tools, run shell commands, spawn subprocesses, and open network connections~\cite{yao2023react}, often without a human reviewing each step. In network operations, postmortems and compliance already require that agent decisions be auditable and replayable~\cite{martinezpaez2026vetra}. Agentic systems also widen the attack surface through prompt injection, tool misuse, and compromised components~\cite{chhabra2026agentic,ma2026openclaw}. In a red-teaming study, tool-enabled agents leaked sensitive data and ran destructive system actions~\cite{shapira2026chaos}. In a reported cyber-espionage campaign, an AI agent carried out most of the intrusion autonomously~\cite{anthropic2025espionage,dawson2026agentic}. Attributing such actions to the execution that performed them is a precondition for accountability, auditing, and containment. When agents run as shared services (\eg a Model Context Protocol (MCP)~\cite{mcp2025spec} server), several executions share a host, a container, or even a worker process. Thus, managing them as network consumers requires answering a simple question: \emph{which agent execution opened this connection?} The answer determines per-execution traffic accounting, the scope of egress policies, and the evidence available when an execution misbehaves.

One straightforward approach to monitoring the actions of such agents is to rely on kernel-level information, such as process creation, file accesses, and network connections, which an agent cannot avoid by switching tools, libraries, or wrappers~\cite{tetragon}. Further, recent technologies, such as the extended Berkeley Packet Filter (eBPF)~\cite{ebpf}, make this practical by attaching verified programs to kernel events, such as process creation and socket operations, without modifying the kernel or the monitored applications~\cite{gbadamosi2024ebpf}. However, the kernel sees connections, processes, and control groups (\textit{cgroups}), but not tasks or tool calls. 

Recent eBPF-based agent systems bridge this gap with different signals: process lineage and temporal correlation~\cite{zheng2025agentsight}, cgroups aligned with tool calls~\cite{zheng2026agentcgroup}, bindings between application context and runtime facts~\cite{byteyellow2026agentprovenance}, inherited task identifiers~\cite{cui2026contractwarden}, and policy domains or tracked process scopes for enforcement~\cite{zheng2026actplane,shukla2026agentguard}. Each signal is informative only under specific execution structures. In a shared \textit{cgroup}, membership cannot separate concurrent executions; inside a shared worker, neither cgroup nor lineage can. Recent work still describes per-tool-call attribution as an open problem~\cite{king2026kernelevidence}.

To address such an open problem, this paper presents \solution{} (Fused Attribution, Accounting, and Control per Task), a framework that attributes each network connection to the agent execution that opened it, reports ambiguity explicitly when the available signals cannot separate executions, and uses the resulting identity for per-execution management. In summary, the contributions of this paper are two-fold: \1 a task-level attribution mechanism that fuses runtime identity, process lineage, and cgroup membership and remains correct for asynchronous tools in shared MCP servers; and \2 per-execution accounting and kernel-enforced egress control on the same identity. Experiments with a real MCP server, a real network interface, and LLM agents driven by open and commercial models show that \solution{} attributes every connection correctly, whereas an existing system misattributes up to 96\% of them, at a cost of microseconds per connection.

The remainder of this paper is organized as follows: Section~\ref{sec:background} provides the background on the involved topics, Section~\ref{sec:related} reviews related work, Section~\ref{sec:design} presents \solution, Section~\ref{sec:method} describes the methodology, Section~\ref{sec:results} reports the results, Section~\ref{sec:discussion} discusses limitations, and Section~\ref{sec:conclusion} concludes this paper.